% !TEX encoding = UTF-8 Unicode
\documentclass[11pt,a4paper]{report}		                                   

\usepackage{graphicx}
\usepackage{a4}
\usepackage[T1]{fontenc}
\usepackage[utf8]{inputenc}
\usepackage[ngerman]{babel}

\usepackage{wrapfig}  
\usepackage{acronym}
\usepackage{enumitem}
\usepackage[ampersand]{easylist}

\usepackage{listings}
\usepackage{color}
\usepackage{graphicx}
\usepackage{float}

\usepackage{caption}
\captionsetup{margin=8pt,font=footnotesize,format=hang}

\usepackage{url}

%Platzierung von Tabelle in Kapitel erzwingen 
\usepackage[section]{placeins}

%Tabellen
\usepackage{booktabs}
\usepackage[table,xcdraw]{xcolor}

\usepackage{chngcntr}
\counterwithout{figure}{chapter} %Kapitelübergreifende Bildnummerierung
\counterwithout{table}{chapter} %Kapitelübergreifende Tabellennummerierung

%Listings
\usepackage{listings} 
\lstset{
	  numbers=none,
	  numberstyle=\tiny,
	  numbersep=5pt,
	  captionpos=b,
	  showstringspaces=false,
	  basicstyle=\footnotesize\ttfamily,
	  numberbychapter=false%TODO
	  }
\lstset{language=Java}
\renewcommand{\lstlistingname}{Auflistung}
\renewcommand{\lstlistlistingname}{Auflistungen}
 
%Baumstruktur
\usepackage{tikz}
\usetikzlibrary{arrows,shapes,positioning,shadows,trees}

\tikzset{
	basic/.style = {draw, text width=2m, drop shadow, font=\sffamily, rectangle},
	root/.style = {basic, rounded corners=2pt, thin, align=center,
		fill=green!30},
	level 2/.style = {basic, rounded corners=6pt, thin,align=center, fill=green!60,
		text width=8em},
	level 3/.style = {basic, thin, align=left, fill=pink!60, text width=6.5em}
}

%Kapitelüberschriften
 \usepackage{fancyhdr} %Fancy Header
 \pagestyle{fancy}
 \fancyhead{}  
 \fancyhead[RE,RO]{\normalsize\sffamily \thechapter \ \leftmark} 
\cfoot{ \thepage}
 \renewcommand{\headrulewidth}{0pt}  
 \renewcommand{\chaptermark}[1]{\markboth{#1}{}}
 
 % set colors for links and URLs
 \usepackage[
 colorlinks=true,
 urlcolor=black,
 linkcolor=black,
 citecolor=black
 ]{hyperref}
 
%Fußnoten
\usepackage{scrextend} %Formatierung der Fuߟnoteneinträge
\deffootnote{1.5em}{1em}{ %Fußnoten Formatierung
	\makebox[1.5em][l]{\thefootnotemark}} % [l]-->Links

%handle subsections
\setcounter{secnumdepth}{3}
\setcounter{tocdepth}{3}

% Literatur - Achtung alphadin file angepasst!
\bibliographystyle{alphadin}




%Eigene Commands, falls gewuenscht
\newcommand{\umbruch}{\\ \noindent}		%Zeilenumbruch
\newcommand{\absatz}{\\ \\ \noindent}	%Absatz ohne einrücken
\newcommand{\extra}[1]{\emph{#1}} 		%Fremdworte, Englische Worte, ... kursiv

\renewcommand\lstlistingname{Quelltext} % Change language of section name

\newcommand{\thema}{Data Analytics}
\newcommand{\untertitel}{Einführung in ESPER}
\newcommand{\autoren}{Lukas Groß, Clemens Grabow}
\newcommand{\abgabedatum}{19.07.2018}
\newcommand{\arbeitsTyp}{Ausarbeitung}

\begin{document}

\begin{titlepage}

\vspace*{-3.5cm}

\begin{flushleft}
\hspace*{-1cm} \includegraphics[width=15.7cm]{front/htwg-logo.png}
\end{flushleft}

\vspace{2.5cm}

\begin{center}
	\huge{
		\textbf{\thema} \\[3cm]
	}
	\Large{
		\textbf{\untertitel}} \\[6.5cm]
	\Large{
		\textbf{\autoren}} \\[3cm]
	\large{
		\textbf{Konstanz, \abgabedatum} \\[1.8cm]
	}
	
	\Huge{
		\textbf{{\sf \arbeitsTyp}}
	}
\end{center}

\end{titlepage}


\pagenumbering{gobble}
\tableofcontents
	\thispagestyle{empty}
\clearpage 
\pagenumbering{arabic}

% So werden Kapitel inkludiert
% \include{chapters/chapterName}
% !TEX encoding = UTF-8 Unicode
\chapter{Motivation}

% !TEX encoding = UTF-8 Unicode
\chapter{Grundlagen}

\section{CEP}
% !TEX encoding = UTF-8 Unicode
\chapter{Konzepte}

\section{Statements}

\section{Patterns}

\section{Filter}

\section{Aggregation}

\section{Data-Windows}

\section{Partitions}
% !TEX encoding = UTF-8 Unicode
\chapter{Architektur}

% !TEX encoding = UTF-8 Unicode
\chapter{Umsetzung}

\section{Installation}

\section{Implementierung}

% !TEX encoding = UTF-8 Unicode
\chapter{Fazit}



\bibliography{bib/literatur_esper}
\addcontentsline{toc}{chapter}{Literaturverzeichnis}

% !TEX encoding = UTF-8 Unicode
\chapter*{Abkürzungsverzeichnis}
\addcontentsline{toc}{chapter}{Abkürzungsverzeichnis}

\begin{acronym}[Bash]
	\acro{Mod}{Modifikation}
	\acro{NPC}{Non-player characters}
\end{acronym}

\end{document}